\subsection{Matched directed opportunities and admitted service}
\label{sec:matched-opportunities}

The matched experiment executes native Continuity with per-direction send
budgets $b\in\{32,64\}$ and the five reference profiles under seven common
schedules. Its 49 configuration--schedule pairs place exposure, admitted
application work and communication cost on the same opportunity budget.
The protocol source is fixed at \texttt{358b921a}; the reference SPQR revision
and its isolated chunk-width variant retain the identities stated above.

\paragraph{Workload and capture.}
Each schedule contains 2,048 directed transmission opportunities, indexed
0--2,047. The designated sender receives one 32-byte application request at
each opportunity and may emit at most one frame of at most 8,192 bytes.
Uncommitted requests remain in its FIFO queue. A response queued by receipt
uses a later opportunity in the response direction. Protocol retries consume
opportunities and retain their exact bytes; network duplication adds only a
delivery. Terminal draining delivers already emitted frames.

Both capture experiments use the cut after opportunity 0's scheduled
local and network actions and its checkpoint, and before opportunity 1
begins local computation. They separately acquire the complete current session
state of $A$ or $B$, together with every public emission, including packets
lost before delivery. For Continuity this is the selected session record and
its matching wrapping key; its independently owned identity signer remains
separate. The references authenticate through their session MAC state.
One application has been committed before the cut in every profile, so the
post-capture denominator is the total first-commitment count minus one.

Native admission is a DeviceJournal-committed ciphertext; reference admission
is a returned component state, with a separately executed 32-byte-payload
AEAD audit. The native predictor derives candidate keys from the captured
epoch-zero chain seeds and checks them against the actual ciphertext, with
altered-seed rejection controls. Reference predictors recover component
message keys, and the execution audits their envelopes with the actual keys
and wrong-key controls. The displayed derived sets are the outputs of these
specified predictors. The cryptographic recovery results use the eligibility
and observation interfaces of Sections~\ref{sec:security-model}
and~\ref{sec:computation}.

\paragraph{Current-epoch capacity and renewed exchange.}
At this cut, the first application has consumed one of the two directions'
epoch-zero allocation indices, and the captured native records contain no
pending rekey entropy. Each direction admits at most $b$ first allocations
under that epoch. Thus at most $2b-1$ further epoch-zero allocations remain;
exact retransmissions reuse their original allocation. Equality requires both
directions to use their remaining epoch-zero capacity. This elementary
allocation count supplies the exposure envelope tested by the native trace
predictor.

% Generated by artifacts/matched-comparison/analyze_matched.py.
\begin{table}[t]
\centering
\small
\setlength{\tabcolsep}{3.5pt}
\begin{tabular}{@{}lrrrr@{}}
\toprule
Profile & Committed / & Capture $A$ & Capture $B$ & Emitted\\
 & delivered & derived / admitted & derived / admitted & bytes\\
\midrule
Continuity, $b=32$ & 688 / 688 & 63 / 687 & 63 / 687 & 565,288\\
Continuity, $b=64$ & 718 / 718 & 127 / 717 & 127 / 717 & 407,596\\
SPQR, chunk 32 & 2,048 / 2,048 & 172 / 2,047 & 85 / 2,047 & 172,339\\
SPQR, chunk 64 & 2,048 / 2,048 & 88 / 2,047 & 43 / 2,047 & 227,268\\
Whole KEM, interval 1 & 2,048 / 2,048 & 4 / 2,047 & 1 / 2,047 & 1,785,156\\
Whole KEM, interval 32 & 2,048 / 2,048 & 66 / 2,047 & 1 / 2,047 & 196,704\\
Whole KEM, interval 64 & 2,048 / 2,048 & 130 / 2,047 & 1 / 2,047 & 157,744\\
\bottomrule
\end{tabular}
\caption{Alternating transport with 2,048 directed opportunities and 2,048 offered 32-byte requests. Each exposure denominator counts first commitments after the capture cut; $A/B$ identify the separately captured endpoint. Emitted bytes include actual native frames or the reference component with its executed payload envelope. Uncommitted requests equal 2,048 minus the committed count.}
\label{tab:matched-opportunities}
\end{table}

\begin{table}[t]
\centering
\small
\setlength{\tabcolsep}{4pt}
\begin{tabular}{@{}rrrrrrrr@{}}
\toprule
Budget & First & ACK & Control & App. & All exact & Image & \texttt{sync\_data}\\
 & app. & & & retries & retries & revisions & calls\\
\midrule
32 & 688 & 676 & 80 & 604 & 636 & 3,554 & 14,219\\
64 & 718 & 711 & 40 & 579 & 596 & 3,644 & 14,579\\
\bottomrule
\end{tabular}
\caption{Native alternating communication and persistence. First application, ACK, control and application-retry columns partition the 2,048 emitted frames. All exact retries also count repeated control frames and overlap that partition. Image revisions and physical sync calls sum both endpoints.}
\label{tab:matched-accounting}
\end{table}


In the alternating schedule, reducing $b$ from 64 to 32 changes the derived
set from 127 of 717 post-capture commitments to 63 of 687
(Table~\ref{tab:matched-opportunities}). Both counts attain $2b-1$.
The same change delivers 688 rather than 718 applications, a reduction of
30, or 4.18\%, under the common 2,048 opportunities. Both endpoints finish
at confirmed epoch 10 for $b=32$ and epoch 5 for $b=64$; the corresponding
emitted-byte costs are 565,288 and 407,596. The budget therefore has a
measured exposure--service--communication tradeoff under this workload.
The references give their own captured-state outcomes at the same cut. For
example, SPQR with 32-byte chunks derives 172 of 2,047 later commitments
from capture $A$ and 85 of 2,047 from capture $B$, reflecting the different
pending material present in the two endpoint snapshots.

\paragraph{Communication policy and durability.}
The native driver uses fair cyclic arbitration among queued control traffic,
ACKs, exact application retries and new applications. An unacknowledged frame
is immediately eligible for retry. Table~\ref{tab:matched-accounting} exposes
that policy's cost: the alternating $b=32$ execution spends 604 opportunities
on application retries and the $b=64$ execution spends 579. Independent
trace analysis places 592 and 572 of those transmissions after the request's
first logical delivery, accounting for 83,472 and 80,652 bytes. This uses the
analyst's delivery history; the sender is still awaiting its ACK. The four
frame classes partition each run's 2,048 emissions. Image-revision and
\texttt{sync\_data} counts measure the real native journal. Each alternating
reference execution returns component state 4,096 times; its physical
persistence fields are \texttt{null}, recording the component-serialization
boundary. The communication totals include all native control, ACK and retry
frames, and the references' complete component-plus-payload envelopes.

\paragraph{Transport controls.}
The one-way schedule supplies only $A\!\to\!B$ opportunities, including for
control traffic. Continuity commits and delivers 32 or 64 applications and
then retains 2,016 or 1,984 offered requests in its queue. Both endpoints
remain at confirmed epoch zero; all 31 or 63 post-capture commitments are
derived from either captured state. The references continue committing
2,048 requests, with 2,047 of 2,047 post-capture keys derived. This control
identifies budget-induced admission suspension and its backlog explicitly.
Recovery during renewed exchange is assessed through the later admitted
allocations and their original sources.

The lossy schedule drops the 205 common opportunities whose index is
congruent to 4 modulo 10. Continuity delivers all 665 or 681 admitted
applications, charging retransmissions to the opportunity budget. Each
reference commits 2,048 component applications and delivers 1,843 under its
declared workload without application retransmission. Network duplication
leaves native first commitments, logical deliveries, emitted bytes, final
epochs and derived sets equal to the alternating execution. Reordering, 9:1
asymmetry and a 256-opportunity offline prefix complete the seven-schedule
inventory; their request queues, exact derived sets and work counters are
retained in the complete 49-row artifact.

The manuscript tables are regenerated from the hash-bound executed report
and checked against its independently recomputed public-evidence projection.
Appendix~\ref{app:matched-comparison} identifies the source, execution and
validation records.
